% Generated draft: requires mathematical and source review.
\documentclass[11pt]{article}
\usepackage[margin=25mm]{geometry}
\usepackage{fontspec}
\setmainfont{Latin Modern Roman}
\usepackage{amsmath,amssymb,mathtools,unicode-math}
\setmathfont{Latin Modern Math}
\usepackage{enumitem,needspace,xurl,hyperref}
\setlength{\parindent}{0pt}
\setlength{\parskip}{5pt}
\setlength{\emergencystretch}{3em}
\newcommand{\N}{\mathbb{N}}
\newcommand{\Z}{\mathbb{Z}}
\newcommand{\Q}{\mathbb{Q}}
\newcommand{\R}{\mathbb{R}}
\newcommand{\C}{\mathbb{C}}
\newcommand{\F}{\mathbb{F}}
\newcommand{\A}{\mathbb{A}}
\newcommand{\PP}{\mathbb{P}}
\newcommand{\E}{\mathbb{E}}
\newcommand{\eps}{\varepsilon}
\newcommand{\dd}{\,\mathrm{d}}
\begin{document}
\section*{UnsolvedMath ID: 30002048}
\textbf{Source problem code:} catalogue entry
\subsection{Definitions and mathematical statement}
Let \(d\geq 7\) be an integer and let \(\alpha\) be an algebraic integer with \([\Q(\alpha):\Q]=d\). Put \(K=\Q(\alpha)\), and let \(\mathcal O_K\) denote its ring of integers. Interpret a unit as an element of \(\mathcal O_K^{\times}\), that is, an element having a multiplicative inverse in \(\mathcal O_K\).

The source defines \(E_0(\alpha)=0\) if \(\alpha-1\) is not a unit; otherwise it defines
\[
E_0(\alpha)=\max\bigl\{n\in\Z_{\geq 0}:\ \alpha^j-1\in\mathcal O_K^{\times}\text{ for every integer }1\leq j\leq n\bigr\},
\]
provided this largest integer exists. It then writes
\[
e(d)=\max\bigl\{E_0(\alpha):\ \alpha\text{ is an algebraic integer and }[\Q(\alpha):\Q]=d\bigr\}.
\]
The supplied conjecture is \(e(d)<d\) for every integer \(d\geq 7\).

A precise formulation of the intended bound, without presupposing existence of these maxima, is:
\[
\begin{gathered}
\text{For every integer }d\geq 7\text{ and every algebraic integer }\alpha
\text{ with }[\Q(\alpha):\Q]=d,\\
\text{there exists an integer }j\text{ with }1\leq j\leq d
\text{ such that }\alpha^j-1\notin\mathcal O_{\Q(\alpha)}^{\times}.
\end{gathered}
\]
This formulation expresses the supplied inequality whenever the source's maxima are defined; if true, it also ensures that those maxima exist in the stated degree range. The supplied record does not explain their existence separately. The ring used for units is the standard number-field interpretation stated above; the source excerpt does not specify it explicitly.
\subsection{Short English statement}
The conjecture asks whether every algebraic integer of degree \(d\geq 7\) has at least one nonunit among \(\alpha-1,\alpha^2-1,\ldots,\alpha^d-1\). The supplied record labels this problem open; that status and its source-verification claim have not been independently checked.
\subsection{Sources}
\begin{itemize}
\item Supplied local UnsolvedMath registry record: permanent ID \(30002048\), title ``Exceptional-Unit Bounds by Algebraic Degree,'' source problem code \texttt{catalogue entry}, in the Oberwolfach Reports Open Problems collection. No individual proposer is supplied; the recorded proposed year is 2012.
\item Supplied report citation: \emph{Diophantische Approximationen (2013)}, Oberwolfach Reports, DOI 10.4171/owr/2012/22, pp. 1307--1359. \url{https://doi.org/10.4171/owr/2012/22}. The record describes Oberwolfach Reports as documenting mathematical workshops and their open-problem discussions. Its statement status is \texttt{corrected\_verified}, with the supplied claim that the official report gives this definition and the same degree threshold; this claim has not been independently verified here.
\item Hare--Mossinghoff, \emph{Exceptional units and cyclic resultants}. \url{https://www.impan.pl/shop/publication/transaction/download/product/82792}.
\item Hare--Mossinghoff, \emph{Exceptional units and cyclic resultants II}. \url{https://uwaterloo.ca/pure-mathematics/sites/default/files/uploads/documents/exceptional_units-5_0.pdf}.
\item Literature assessment dated exactly 2026-08-21: supplied dated curation, not independently verified. The supplied assessment says that the exceptional-unit literature still records \(e(d)<d\) for \(d\geq 7\) as a conjecture, and that known small-degree computations and lower bounds do not establish it in all degrees. The background characterizes this as dated literature triage derived from the supplied OWR curation and requests independent verification of time-sensitive status and cited claims.
\end{itemize}
\end{document}
